% Seção: F-RES · Falhas e resiliência (RF-RES-01 a RF-RES-03)
% Dono: Caio · Números: docs/plano-de-desenvolvimento/features/feature-F-RES.md, Resultados
% (comandos: python benchmarks/medir_falhas.py, execuções de 10/11 e 19/11/2026;
% python benchmarks/variancia_falhas.py --exata, 19/11/2026). Figura em F-MET.tex (dupla).
% Condensada em 20/11 para o orçamento de 15 páginas.
\subsection{Falhas e resiliência}\label{sec:F-RES}

A ferramenta simula a saída de operação de uma parada, de um trecho ou de uma linha sobre $G_P$
e compara a rede antes e depois. A falha não reconstrói o grafo: desativa os vértices e as
arestas atingidos e os reativa ao fim. Uma linha remove apenas os trechos que nenhuma outra
linha percorre. As componentes fortemente conexas saem do Tarjan iterativo, em $O(V + E)$: a
versão recursiva não serve, porque a busca em profundidade chega a 4.657 níveis no $G_P$ real,
acima do limite de recursão do Python (1.000). A perda de eficiência é
$(E_{\text{antes}} - E_{\text{depois}})/E_{\text{antes}}$, mantido o $n$ de antes da falha, para
que remover uma parada nunca aumente a eficiência; com $s$ origens, custa $O(s(V + E)\log V)$.

No feed de 11/11 ($V = 9{.}659$, $E = 122{.}732$), $G_P$ tem 4 componentes fortemente conexas
(uma com 9.656 paradas e 3 unitárias), iguais às do NetworkX. Com 100 origens, retirar a parada
com mais saídas de viagem custou 0,036\% de eficiência e a linha com mais trechos exclusivos,
0,071\%, sem isolar paradas. Das 291 linhas, retiradas uma por vez, só 17 tiram alguma parada da
maior componente; a pior, a 832, deixa 33 paradas fora (Figura~\ref{fig:falhas}) e custa
0,50\% exatos (todas as 9.659 origens, 22~min). Cada simulação com 100 origens levou de 18 a
38~s (Intel Core i5-1235U); a página usa 30. \textbf{Limitações.} A perda mede a rede
estrutural, não o tempo porta a porta. Quando a falha desconecta poucas paradas, a
perda amostrada tem cauda pesada: na linha 832, cinco sementes com 100 origens deram de 0,24\% a
1,76\%, e a página avisa quando isso pode acontecer.
